% !TeX root = ../../Book3.tex
% 纲要标题：### 14.4 域扩张、可分性与微分
% 纲要位置：工作记录/计划纲要.md，第 2399 行。
% 纲要细目（写作范围）：
% * 单个代数元的微分与有限域扩张的可分性判据
% * 指数一纯不可分扩张；子域 \(K L^p\) 及微分基的构造
% * 纯不可分扩张的次数与微分维数；基变换后的幂零结构

\section{域扩张、可分性与微分}
\label{b3:sec:1404}

代数元的关系 \(f(\alpha)=0\) 对导子给出 \(f'(\alpha)\mathrm{d}\alpha=0\)。若导数在域中非零，这个方向被迫消失；若关系是 \(p\) 次幂，求导却不能看见它。本节从单个代数元的计算出发，说明这种差别如何检测一般有限域扩张的可分性，再比较纯不可分扩张的次数与微分维数。

\subsection{有限代数域扩张的微分模}
\label{b3:subsec:1404:01}

先看一个代数元。设 \(K\) 为域，\(\alpha\) 为某个扩域中对 \(K\) 代数的元素，\(L=K(\alpha)\)，\(f\in K[T]\) 为 \(\alpha\) 的首一不可约多项式。由于 \(L=K[T]/(f)\)，\cref{b3:cor:differential-presentation}给出
\begin{equation}\label{b3:eq:simple-field-differentials}
\Omega_{L/K}\cong L\,\mathrm{d}\alpha/\bigl(f'(\alpha)\mathrm{d}\alpha\bigr).
\end{equation}
若 \(f\) 可分，\(f'(\alpha)\ne0\) 在域中为单位，故微分模为零。若 \(f\) 不可分，其不可约性迫使 \(f'=0\)，于是微分模为一维。一般有限扩张不一定由一个元素生成；不能把这个单生成元计算直接当作一般情形的证明。

\begin{lemma}\label{b3:lem:exponent-one-differentials}
设 \(K\subseteq L\) 为特征 \(p>0\) 的有限域扩张，令 \(F=K L^p\) 为由 \(K\) 及所有 \(p\) 次幂生成的子域。存在 \(\alpha_1,\ldots,\alpha_r\in L\)，使
\[
[L:F]=p^r,\qquad
\Omega_{L/K}\cong\bigoplus_{i=1}^r L\,\mathrm{d}\alpha_i.
\]
其中 \(r=0\) 恰好表示 \(L=K L^p\)。
\end{lemma}
\begin{proof}
每个 \(K\) 导子都零化 \(L^p\)，再由四则运算法则零化 \(F\)，所以 \(K\) 导子与 \(F\) 导子相同，微分模也相同。逐次选取尚未落入已生成子域的 \(\alpha_i\)，直至生成 \(L/F\)。记前一步生成的子域为 \(F_{i-1}=F(\alpha_1,\ldots,\alpha_{i-1})\)。因为 \(\alpha_i^p\in F\)，其在 \(F_{i-1}\) 上的最小多项式整除 \(T^p-\alpha_i^p\)。若次数为 \(j<p\)，则在代数闭包中必为
\[
(T-\alpha_i)^j=T^j-j\alpha_iT^{j-1}+\cdots,
\qquad 1\le j<p.
\]
系数 \(-j\alpha_i\) 属于 \(F_{i-1}\)，而 \(j\) 在特征 \(p\) 的域中非零，故 \(\alpha_i\in F_{i-1}\)，与选取方式矛盾。因此最小多项式恰为 \(T^p-\alpha_i^p\)，每步次数为 \(p\)，而有限次数保证过程终止。

于是 \(\prod_i\alpha_i^{e_i}\)、\(0\le e_i<p\)，组成 \(F\) 基。自然满射
\[
F[X_1,\ldots,X_r]/(X_1^p-\alpha_1^p,\ldots,X_r^p-\alpha_r^p)
\longrightarrow L
\]
两侧维数都是 \(p^r\)，故为同构。关系的微分全部为零，\cref{b3:cor:differential-presentation}遂给出所列自由基。
\end{proof}

\begin{theorem}\label{b3:thm:finite-field-differentials-separable}
设 \(L/K\) 为有限域扩张。则
\[
\Omega_{L/K}=0\quad\Longleftrightarrow\quad L/K\text{ 可分}.
\]
若 \(L/K\) 为特征 \(p>0\) 的非平凡有限纯不可分扩张，则 \(\Omega_{L/K}\ne0\)。
\end{theorem}
\begin{proof}
若扩张可分，对任意 \(\alpha\in L\)，其最小多项式 \(f\) 满足 \(f'(\alpha)\ne0\)。微分等式 \(f'(\alpha)\mathrm{d}\alpha=0\) 因而给出 \(\mathrm{d}\alpha=0\)。这些元素生成微分模，所以该模为零。

先证非平凡纯不可分情形。设 \(L/K\) 的指数不超过 \(e\)，即 \(L^{p^e}\subseteq K\)；有限生成性保证存在同一个 \(e\)。若 \(L=K L^p\)，反复取 \(p\) 次幂并代回便得
\(L=K L^{p^j}\) 对所有 \(j\ge1\) 成立，特别地 \(L=K\)，矛盾。因此 \(K L^p\subsetneq L\)，前一引理给出至少一个非零微分。

最后设 \(L/K\) 不可分。令 \(E\) 为 \(L\) 中对 \(K\) 可分的元素组成的子域；可分代数扩张对复合及子扩张封闭，保证这些元素对四则运算封闭。\(E/K\) 有限可分。对任意 \(\alpha\in L\)，其不可约多项式在特征 \(p\) 下唯一写成 \(g(T^{p^a})\)，其中 \(g'\ne0\)。多项式 \(g\) 仍不可约，否则将其分解代入 \(T^{p^a}\) 就得到原多项式的分解；故 \(g\) 可分，\(\alpha^{p^a}\) 对 \(K\) 可分并属于 \(E\)。于是 \(L/E\) 纯不可分；由于原扩张不可分，它非平凡。复合微分正合列给出
\[
L\otimes_E\Omega_{E/K}\longrightarrow\Omega_{L/K}
\longrightarrow\Omega_{L/E}\longrightarrow0.
\]
左端由已证的可分情形为零，右端由刚证的纯不可分情形非零，所以中间项非零。特征零没有不可分代数扩张，不需此分情形。
\end{proof}

\subsection{纯不可分扩张的微分计算}
\label{b3:subsec:1404:05}

微分维数记录独立的一阶方向，却未必记录纯不可分扩张的全部幂次深度。设 \(p\) 为素数，\(s\) 为超越元，\(k=\mathbb F_p(s)\)、\(L=k(\alpha)\)，其中 \(\alpha^{p^e}=s\)、\(e\ge1\) 为整数。则
\[
\Omega_{L/k}=L\,\mathrm{d}\alpha.
\]
最小多项式 \(T^{p^e}-s\) 的不可约性可用 \(s\) 在 \(\mathbb F_p[s]\) 中的 Eisenstein 判据检验。扩张次数为 \(p^e\)，微分维数却一直为一，因此微分维数不是不可分次数的对数。它测量的是独立的一阶方向，未记录每个方向的全部幂次深度。

\ProseParagraphBreak

多方向的例子可由两个代数独立的元 \(s,t\) 构造。设 \(p\) 为素数，取
\[
k=\mathbb F_p(s,t),\qquad
L=k(\alpha,\beta),\qquad \alpha^p=s,\quad\beta^p=t.
\]
有 \([L:k]=p^2\) 及
\(\Omega_{L/k}=L\,\mathrm{d}\alpha\oplus L\,\mathrm{d}\beta\)。这里可把 \(L\) 识别为有理函数域 \(\mathbb F_p(\alpha,\beta)\)：不同模 \(p\) 的指数类在多项式恒等式中不会相消，所以 \(\alpha^i\beta^j\)、\(0\le i,j<p\)，对 \(k=\mathbb F_p(\alpha^p,\beta^p)\) 线性无关，并显然生成。对应的两个关系微分都为零，给出所列基。

\ProseParagraphBreak
在这一例子中，原来的 \(L\) 是域，扩到 \(L\) 本身后却有
\[
L\otimes_kL\cong L[\varepsilon,\eta]/(\varepsilon^p,\eta^p).
\]
把右边变量分别送到第二因子的 \(\alpha,\beta\) 减去第一因子的相应元素，便得到满射；两侧 \(L\) 维数都是 \(p^2\)，所以为同构。与\cref{b3:exm:inseparable-tangent}一样，这个计算把不可分性表现为基变换后出现的一阶方向和幂零结构。
